\documentclass[
  doc,
  floatsintext,
  longtable,
  a4paper,
  nolmodern,
  notxfonts,
  notimes,
  12pt,
  colorlinks=true,linkcolor=blue,citecolor=blue,urlcolor=blue]{apa7}

\usepackage{amsmath}
\usepackage{amssymb}

\geometry{inner=1in, outer=1in}
\fancyhfoffset[LE,RO]{0cm}


\usepackage[bidi=default]{babel}
\babelprovide[main,import]{spanish}


% get rid of language-specific shorthands (see #6817):
\let\LanguageShortHands\languageshorthands
\def\languageshorthands#1{}

\RequirePackage{longtable}
\RequirePackage{threeparttablex}

\makeatletter
\renewcommand{\paragraph}{\@startsection{paragraph}{4}{\parindent}%
	{0\baselineskip \@plus 0.2ex \@minus 0.2ex}%
	{-.5em}%
	{\normalfont\normalsize\bfseries\typesectitle}}

\renewcommand{\subparagraph}[1]{\@startsection{subparagraph}{5}{0.5em}%
	{0\baselineskip \@plus 0.2ex \@minus 0.2ex}%
	{-\z@\relax}%
	{\normalfont\normalsize\bfseries\itshape\hspace{\parindent}{#1}\textit{\addperi}}{\relax}}
\makeatother




\usepackage{longtable, booktabs, multirow, multicol, colortbl, hhline, caption, array, float, xpatch}
\usepackage{subcaption}


\renewcommand\thesubfigure{\Alph{subfigure}}
\setcounter{topnumber}{2}
\setcounter{bottomnumber}{2}
\setcounter{totalnumber}{4}
\renewcommand{\topfraction}{0.85}
\renewcommand{\bottomfraction}{0.85}
\renewcommand{\textfraction}{0.15}
\renewcommand{\floatpagefraction}{0.7}

\usepackage{tcolorbox}
\tcbuselibrary{listings,theorems, breakable, skins}
\usepackage{fontawesome5}

\definecolor{quarto-callout-color}{HTML}{909090}
\definecolor{quarto-callout-note-color}{HTML}{0758E5}
\definecolor{quarto-callout-important-color}{HTML}{CC1914}
\definecolor{quarto-callout-warning-color}{HTML}{EB9113}
\definecolor{quarto-callout-tip-color}{HTML}{00A047}
\definecolor{quarto-callout-caution-color}{HTML}{FC5300}
\definecolor{quarto-callout-color-frame}{HTML}{ACACAC}
\definecolor{quarto-callout-note-color-frame}{HTML}{4582EC}
\definecolor{quarto-callout-important-color-frame}{HTML}{D9534F}
\definecolor{quarto-callout-warning-color-frame}{HTML}{F0AD4E}
\definecolor{quarto-callout-tip-color-frame}{HTML}{02B875}
\definecolor{quarto-callout-caution-color-frame}{HTML}{FD7E14}

%\newlength\Oldarrayrulewidth
%\newlength\Oldtabcolsep


\usepackage{hyperref}



\usepackage{color}
\usepackage{fancyvrb}
\newcommand{\VerbBar}{|}
\newcommand{\VERB}{\Verb[commandchars=\\\{\}]}
\DefineVerbatimEnvironment{Highlighting}{Verbatim}{commandchars=\\\{\}}
% Add ',fontsize=\small' for more characters per line
\usepackage{framed}
\definecolor{shadecolor}{RGB}{241,243,245}
\newenvironment{Shaded}{\begin{snugshade}}{\end{snugshade}}
\newcommand{\AlertTok}[1]{\textcolor[rgb]{0.68,0.00,0.00}{#1}}
\newcommand{\AnnotationTok}[1]{\textcolor[rgb]{0.37,0.37,0.37}{#1}}
\newcommand{\AttributeTok}[1]{\textcolor[rgb]{0.40,0.45,0.13}{#1}}
\newcommand{\BaseNTok}[1]{\textcolor[rgb]{0.68,0.00,0.00}{#1}}
\newcommand{\BuiltInTok}[1]{\textcolor[rgb]{0.00,0.23,0.31}{#1}}
\newcommand{\CharTok}[1]{\textcolor[rgb]{0.13,0.47,0.30}{#1}}
\newcommand{\CommentTok}[1]{\textcolor[rgb]{0.37,0.37,0.37}{#1}}
\newcommand{\CommentVarTok}[1]{\textcolor[rgb]{0.37,0.37,0.37}{\textit{#1}}}
\newcommand{\ConstantTok}[1]{\textcolor[rgb]{0.56,0.35,0.01}{#1}}
\newcommand{\ControlFlowTok}[1]{\textcolor[rgb]{0.00,0.23,0.31}{\textbf{#1}}}
\newcommand{\DataTypeTok}[1]{\textcolor[rgb]{0.68,0.00,0.00}{#1}}
\newcommand{\DecValTok}[1]{\textcolor[rgb]{0.68,0.00,0.00}{#1}}
\newcommand{\DocumentationTok}[1]{\textcolor[rgb]{0.37,0.37,0.37}{\textit{#1}}}
\newcommand{\ErrorTok}[1]{\textcolor[rgb]{0.68,0.00,0.00}{#1}}
\newcommand{\ExtensionTok}[1]{\textcolor[rgb]{0.00,0.23,0.31}{#1}}
\newcommand{\FloatTok}[1]{\textcolor[rgb]{0.68,0.00,0.00}{#1}}
\newcommand{\FunctionTok}[1]{\textcolor[rgb]{0.28,0.35,0.67}{#1}}
\newcommand{\ImportTok}[1]{\textcolor[rgb]{0.00,0.46,0.62}{#1}}
\newcommand{\InformationTok}[1]{\textcolor[rgb]{0.37,0.37,0.37}{#1}}
\newcommand{\KeywordTok}[1]{\textcolor[rgb]{0.00,0.23,0.31}{\textbf{#1}}}
\newcommand{\NormalTok}[1]{\textcolor[rgb]{0.00,0.23,0.31}{#1}}
\newcommand{\OperatorTok}[1]{\textcolor[rgb]{0.37,0.37,0.37}{#1}}
\newcommand{\OtherTok}[1]{\textcolor[rgb]{0.00,0.23,0.31}{#1}}
\newcommand{\PreprocessorTok}[1]{\textcolor[rgb]{0.68,0.00,0.00}{#1}}
\newcommand{\RegionMarkerTok}[1]{\textcolor[rgb]{0.00,0.23,0.31}{#1}}
\newcommand{\SpecialCharTok}[1]{\textcolor[rgb]{0.37,0.37,0.37}{#1}}
\newcommand{\SpecialStringTok}[1]{\textcolor[rgb]{0.13,0.47,0.30}{#1}}
\newcommand{\StringTok}[1]{\textcolor[rgb]{0.13,0.47,0.30}{#1}}
\newcommand{\VariableTok}[1]{\textcolor[rgb]{0.07,0.07,0.07}{#1}}
\newcommand{\VerbatimStringTok}[1]{\textcolor[rgb]{0.13,0.47,0.30}{#1}}
\newcommand{\WarningTok}[1]{\textcolor[rgb]{0.37,0.37,0.37}{\textit{#1}}}

\providecommand{\tightlist}{%
  \setlength{\itemsep}{0pt}\setlength{\parskip}{0pt}}
\usepackage{longtable,booktabs,array}
\newcounter{none} % for unnumbered tables
\usepackage{calc} % for calculating minipage widths
% Correct order of tables after \paragraph or \subparagraph
\usepackage{etoolbox}
\makeatletter
\patchcmd\longtable{\par}{\if@noskipsec\mbox{}\fi\par}{}{}
\makeatother
% Allow footnotes in longtable head/foot
\IfFileExists{footnotehyper.sty}{\usepackage{footnotehyper}}{\usepackage{footnote}}
\makesavenoteenv{longtable}

\usepackage{graphicx}
\makeatletter
\newsavebox\pandoc@box
\newcommand*\pandocbounded[1]{% scales image to fit in text height/width
  \sbox\pandoc@box{#1}%
  \Gscale@div\@tempa{\textheight}{\dimexpr\ht\pandoc@box+\dp\pandoc@box\relax}%
  \Gscale@div\@tempb{\linewidth}{\wd\pandoc@box}%
  \ifdim\@tempb\p@<\@tempa\p@\let\@tempa\@tempb\fi% select the smaller of both
  \ifdim\@tempa\p@<\p@\scalebox{\@tempa}{\usebox\pandoc@box}%
  \else\usebox{\pandoc@box}%
  \fi%
}
% Set default figure placement to htbp
\def\fps@figure{htbp}
\makeatother







\usepackage{newtx}

\defaultfontfeatures{Scale=MatchLowercase}
\defaultfontfeatures[\rmfamily]{Ligatures=TeX,Scale=1}





\title{Conversión de tipos (casting)}


\shorttitle{CONVERSIÓN DE TIPOS (CASTING)}


\usepackage{etoolbox}



\ccoppy{\textcopyright~2026}



\author{Edison Achalma}



\affiliation{
{Escuela Profesional de Economía, Universidad Nacional de San Cristóbal
de Huamanga}}




\leftheader{Achalma}

\date{2026-10-20}


\abstract{Este abstract será actualizado una vez que se complete el
contenido final del artículo. }



\authornote{\par{\addORCIDlink{Edison Achalma}{0000-0001-6996-3364}} 
\par{ }
\par{   El autor no tiene conflictos de interés que revelar.    Los
roles de autor se clasificaron utilizando la taxonomía de roles de
colaborador (CRediT; https://credit.niso.org/) de la siguiente
manera:  Edison Achalma:   conceptualización, metodología, análisis
formal, investigación, recursos, curación de
datos, redacción, visualización, supervisión, administración del
proyecto}
\par{La correspondencia relativa a este artículo debe dirigirse a Edison
Achalma, Escuela Profesional de Economía, Universidad Nacional de San
Cristóbal de Huamanga, Portal Independencia N°
57, Ayacucho, AYA 5001, Perú, Email: \href{mailto:elmer.achalma.09@unsch.edu.pe}{elmer.achalma.09@unsch.edu.pe}}
}

\makeatletter
\let\endoldlt\endlongtable
\def\endlongtable{
\hline
\endoldlt
}
\makeatother

\urlstyle{same}



\hypersetup{
  pdftitle={Conversi\'{o}n de tipos (casting)},
  pdfauthor={Autor}
}
\makeatletter
\@ifpackageloaded{caption}{}{\usepackage{caption}}
\AtBeginDocument{%
\ifdefined\contentsname
  \renewcommand*\contentsname{Tabla de contenidos}
\else
  \newcommand\contentsname{Tabla de contenidos}
\fi
\ifdefined\listfigurename
  \renewcommand*\listfigurename{Lista de Figuras}
\else
  \newcommand\listfigurename{Lista de Figuras}
\fi
\ifdefined\listtablename
  \renewcommand*\listtablename{Lista de Tablas}
\else
  \newcommand\listtablename{Lista de Tablas}
\fi
\ifdefined\figurename
  \renewcommand*\figurename{Figura}
\else
  \newcommand\figurename{Figura}
\fi
\ifdefined\tablename
  \renewcommand*\tablename{Tabla}
\else
  \newcommand\tablename{Tabla}
\fi
}
\@ifpackageloaded{float}{}{\usepackage{float}}
\floatstyle{ruled}
\@ifundefined{c@chapter}{\newfloat{codelisting}{h}{lop}}{\newfloat{codelisting}{h}{lop}[chapter]}
\floatname{codelisting}{Listado}
\newcommand*\listoflistings{\listof{codelisting}{Lista de Listados}}
\makeatother
\makeatletter
\makeatother
\makeatletter
\@ifpackageloaded{caption}{}{\usepackage{caption}}
\@ifpackageloaded{subcaption}{}{\usepackage{subcaption}}
\makeatother
\makeatletter
\@ifpackageloaded{fontawesome5}{}{\usepackage{fontawesome5}}
\makeatother

% From https://tex.stackexchange.com/a/645996/211326
%%% apa7 doesn't want to add appendix section titles in the toc
%%% let's make it do it
\makeatletter
\xpatchcmd{\appendix}
  {\par}
  {\addcontentsline{toc}{section}{\@currentlabelname}\par}
  {}{}
\makeatother

%% Disable longtable counter
%% https://tex.stackexchange.com/a/248395/211326

\usepackage{etoolbox}

\makeatletter
\patchcmd{\LT@caption}
  {\bgroup}
  {\bgroup\global\LTpatch@captiontrue}
  {}{}
\patchcmd{\longtable}
  {\par}
  {\par\global\LTpatch@captionfalse}
  {}{}
\apptocmd{\endlongtable}
  {\ifLTpatch@caption\else\addtocounter{table}{-1}\fi}
  {}{}
\newif\ifLTpatch@caption
\makeatother

\begin{document}

\maketitle


\hypertarget{toc}{}
\tableofcontents
\newpage
\section[Introduction]{Conversión de tipos (casting)}

\setcounter{secnumdepth}{3}

\setlength\LTleft{0pt}




\subsection{🎯 Motivación}\label{motivaciuxf3n}

En programación, a menudo necesitas convertir datos de un tipo a otro.
Por ejemplo, cuando recibes entrada del usuario (siempre es texto), pero
necesitas hacer cálculos matemáticos. O cuando quieres mostrar un número
como parte de un mensaje de texto. El casting te permite controlar
explícitamente qué tipo de dato estás usando, evitando errores y
asegurando que tu código funcione correctamente.

\subsection{💡 Intuición}\label{intuiciuxf3n}

Imagina que tienes diferentes moldes para galletas: uno para círculos
(int), uno para estrellas (float), y otro para letras (str). Tienes masa
(datos) y puedes usar un molde para darle la forma que necesitas. El
casting es como elegir el molde adecuado: \texttt{int()} convierte en
entero, \texttt{float()} en decimal, \texttt{str()} en texto. Pero ten
cuidado: no todos los moldes funcionan con todos los tipos de masa.

\subsection{🧠 Modelo mental}\label{modelo-mental}

Python usa \textbf{funciones constructoras} para crear objetos de tipos
específicos:

\begin{enumerate}
\def\labelenumi{\arabic{enumi}.}
\tightlist
\item
  \textbf{\texttt{int()}}: convierte a entero (elimina decimales,
  requiere texto numérico).
\item
  \textbf{\texttt{float()}}: convierte a decimal (puede convertir
  enteros y texto numérico).
\item
  \textbf{\texttt{str()}}: convierte a texto (funciona con casi
  cualquier cosa).
\item
  \textbf{\texttt{bool()}}: convierte a booleano (reglas específicas de
  verdadero/falso).
\end{enumerate}

El casting es explícito: tú decides cuándo y cómo convertir.

\subsection{📐 Desarrollo}\label{desarrollo}

\subsubsection{\texorpdfstring{1. \texttt{int()} - Convertir a
Entero}{1. int() - Convertir a Entero}}\label{int---convertir-a-entero}

Convierte números (enteros, decimales) y cadenas de texto numéricas a
enteros:

\begin{Shaded}
\begin{Highlighting}[]
\CommentTok{\# Desde enteros}
\NormalTok{x }\OperatorTok{=} \BuiltInTok{int}\NormalTok{(}\DecValTok{1}\NormalTok{)      }\CommentTok{\# 1}
\NormalTok{y }\OperatorTok{=} \BuiltInTok{int}\NormalTok{(}\OperatorTok{{-}}\DecValTok{5}\NormalTok{)     }\CommentTok{\# {-}5}

\CommentTok{\# Desde floats (trunca, no redondea)}
\NormalTok{z }\OperatorTok{=} \BuiltInTok{int}\NormalTok{(}\FloatTok{2.8}\NormalTok{)    }\CommentTok{\# 2 (pierde el .8)}
\NormalTok{w }\OperatorTok{=} \BuiltInTok{int}\NormalTok{(}\OperatorTok{{-}}\FloatTok{3.9}\NormalTok{)   }\CommentTok{\# {-}3 (trunca hacia cero)}

\CommentTok{\# Desde strings (solo si representan números enteros)}
\NormalTok{a }\OperatorTok{=} \BuiltInTok{int}\NormalTok{(}\StringTok{"10"}\NormalTok{)   }\CommentTok{\# 10}
\NormalTok{b }\OperatorTok{=} \BuiltInTok{int}\NormalTok{(}\StringTok{"{-}5"}\NormalTok{)   }\CommentTok{\# {-}5}

\CommentTok{\# ❌ Error: no se puede convertir texto no numérico}
\CommentTok{\# c = int("Hola")  \# ValueError}
\end{Highlighting}
\end{Shaded}

\subsubsection{\texorpdfstring{2. \texttt{float()} - Convertir a
Decimal}{2. float() - Convertir a Decimal}}\label{float---convertir-a-decimal}

Convierte enteros, otros decimales y cadenas numéricas a floats:

\begin{Shaded}
\begin{Highlighting}[]
\CommentTok{\# Desde enteros}
\NormalTok{x }\OperatorTok{=} \BuiltInTok{float}\NormalTok{(}\DecValTok{1}\NormalTok{)     }\CommentTok{\# 1.0}
\NormalTok{y }\OperatorTok{=} \BuiltInTok{float}\NormalTok{(}\OperatorTok{{-}}\DecValTok{5}\NormalTok{)    }\CommentTok{\# {-}5.0}

\CommentTok{\# Desde floats (no cambia)}
\NormalTok{z }\OperatorTok{=} \BuiltInTok{float}\NormalTok{(}\FloatTok{2.8}\NormalTok{)   }\CommentTok{\# 2.8}

\CommentTok{\# Desde strings numéricos}
\NormalTok{a }\OperatorTok{=} \BuiltInTok{float}\NormalTok{(}\StringTok{"3"}\NormalTok{)   }\CommentTok{\# 3.0}
\NormalTok{b }\OperatorTok{=} \BuiltInTok{float}\NormalTok{(}\StringTok{"3.14"}\NormalTok{) }\CommentTok{\# 3.14}
\NormalTok{c }\OperatorTok{=} \BuiltInTok{float}\NormalTok{(}\StringTok{"{-}5.5"}\NormalTok{) }\CommentTok{\# {-}5.5}
\NormalTok{d }\OperatorTok{=} \BuiltInTok{float}\NormalTok{(}\StringTok{"1e3"}\NormalTok{) }\CommentTok{\# 1000.0 (notación científica)}

\CommentTok{\# ❌ Error: texto no numérico}
\CommentTok{\# e = float("Hola")  \# ValueError}
\end{Highlighting}
\end{Shaded}

\subsubsection{\texorpdfstring{3. \texttt{str()} - Convertir a
Texto}{3. str() - Convertir a Texto}}\label{str---convertir-a-texto}

Convierte prácticamente cualquier cosa a texto:

\begin{Shaded}
\begin{Highlighting}[]
\CommentTok{\# Desde enteros}
\NormalTok{x }\OperatorTok{=} \BuiltInTok{str}\NormalTok{(}\DecValTok{5}\NormalTok{)       }\CommentTok{\# "5"}
\NormalTok{y }\OperatorTok{=} \BuiltInTok{str}\NormalTok{(}\OperatorTok{{-}}\DecValTok{3}\NormalTok{)      }\CommentTok{\# "{-}3"}

\CommentTok{\# Desde floats}
\NormalTok{z }\OperatorTok{=} \BuiltInTok{str}\NormalTok{(}\FloatTok{3.14}\NormalTok{)    }\CommentTok{\# "3.14"}
\NormalTok{w }\OperatorTok{=} \BuiltInTok{str}\NormalTok{(}\FloatTok{2.0}\NormalTok{)     }\CommentTok{\# "2.0"}

\CommentTok{\# Desde booleanos}
\NormalTok{a }\OperatorTok{=} \BuiltInTok{str}\NormalTok{(}\VariableTok{True}\NormalTok{)    }\CommentTok{\# "True"}
\NormalTok{b }\OperatorTok{=} \BuiltInTok{str}\NormalTok{(}\VariableTok{False}\NormalTok{)   }\CommentTok{\# "False"}

\CommentTok{\# Desde None}
\NormalTok{c }\OperatorTok{=} \BuiltInTok{str}\NormalTok{(}\VariableTok{None}\NormalTok{)    }\CommentTok{\# "None"}

\CommentTok{\# Desde estructuras de datos}
\NormalTok{d }\OperatorTok{=} \BuiltInTok{str}\NormalTok{([}\DecValTok{1}\NormalTok{, }\DecValTok{2}\NormalTok{, }\DecValTok{3}\NormalTok{])  }\CommentTok{\# "[1, 2, 3]"}
\NormalTok{e }\OperatorTok{=} \BuiltInTok{str}\NormalTok{(\{}\StringTok{"a"}\NormalTok{: }\DecValTok{1}\NormalTok{\})   }\CommentTok{\# "\{\textquotesingle{}a\textquotesingle{}: 1\}"}
\end{Highlighting}
\end{Shaded}

\subsubsection{\texorpdfstring{4. \texttt{bool()} - Convertir a
Booleano}{4. bool() - Convertir a Booleano}}\label{bool---convertir-a-booleano}

Convierte valores a \texttt{True} o \texttt{False} según reglas
específicas:

\begin{Shaded}
\begin{Highlighting}[]
\CommentTok{\# Falsy (valores que son False)}
\BuiltInTok{print}\NormalTok{(}\BuiltInTok{bool}\NormalTok{(}\DecValTok{0}\NormalTok{))      }\CommentTok{\# False}
\BuiltInTok{print}\NormalTok{(}\BuiltInTok{bool}\NormalTok{(}\FloatTok{0.0}\NormalTok{))    }\CommentTok{\# False}
\BuiltInTok{print}\NormalTok{(}\BuiltInTok{bool}\NormalTok{(}\StringTok{""}\NormalTok{))     }\CommentTok{\# False}
\BuiltInTok{print}\NormalTok{(}\BuiltInTok{bool}\NormalTok{([]))     }\CommentTok{\# False}
\BuiltInTok{print}\NormalTok{(}\BuiltInTok{bool}\NormalTok{(\{\}))     }\CommentTok{\# False}
\BuiltInTok{print}\NormalTok{(}\BuiltInTok{bool}\NormalTok{(}\VariableTok{None}\NormalTok{))   }\CommentTok{\# False}

\CommentTok{\# Truthy (valores que son True)}
\BuiltInTok{print}\NormalTok{(}\BuiltInTok{bool}\NormalTok{(}\DecValTok{1}\NormalTok{))      }\CommentTok{\# True}
\BuiltInTok{print}\NormalTok{(}\BuiltInTok{bool}\NormalTok{(}\OperatorTok{{-}}\DecValTok{5}\NormalTok{))     }\CommentTok{\# True}
\BuiltInTok{print}\NormalTok{(}\BuiltInTok{bool}\NormalTok{(}\StringTok{"Hola"}\NormalTok{)) }\CommentTok{\# True}
\BuiltInTok{print}\NormalTok{(}\BuiltInTok{bool}\NormalTok{(}\StringTok{" "}\NormalTok{))    }\CommentTok{\# True (texto con espacio)}
\BuiltInTok{print}\NormalTok{(}\BuiltInTok{bool}\NormalTok{([}\DecValTok{1}\NormalTok{, }\DecValTok{2}\NormalTok{])) }\CommentTok{\# True}
\BuiltInTok{print}\NormalTok{(}\BuiltInTok{bool}\NormalTok{(\{}\StringTok{"a"}\NormalTok{: }\DecValTok{1}\NormalTok{\})) }\CommentTok{\# True}
\end{Highlighting}
\end{Shaded}

\subsubsection{5. Conversión entre Tipos
Numéricos}\label{conversiuxf3n-entre-tipos-numuxe9ricos}

\begin{Shaded}
\begin{Highlighting}[]
\CommentTok{\# Entero → Flotante}
\NormalTok{x }\OperatorTok{=} \DecValTok{5}
\NormalTok{y }\OperatorTok{=} \BuiltInTok{float}\NormalTok{(x)  }\CommentTok{\# 5.0}

\CommentTok{\# Flotante → Entero (trunca)}
\NormalTok{a }\OperatorTok{=} \FloatTok{3.99}
\NormalTok{b }\OperatorTok{=} \BuiltInTok{int}\NormalTok{(a)    }\CommentTok{\# 3 (pierde decimales)}

\CommentTok{\# String numérico → Entero}
\NormalTok{edad }\OperatorTok{=} \StringTok{"25"}
\NormalTok{edad\_int }\OperatorTok{=} \BuiltInTok{int}\NormalTok{(edad)  }\CommentTok{\# 25}

\CommentTok{\# String numérico → Flotante}
\NormalTok{precio }\OperatorTok{=} \StringTok{"19.99"}
\NormalTok{precio\_float }\OperatorTok{=} \BuiltInTok{float}\NormalTok{(precio)  }\CommentTok{\# 19.99}

\CommentTok{\# Entero → String (para mostrar)}
\NormalTok{total }\OperatorTok{=} \DecValTok{100}
\NormalTok{mensaje }\OperatorTok{=} \StringTok{"El total es: "} \OperatorTok{+} \BuiltInTok{str}\NormalTok{(total)  }\CommentTok{\# "El total es: 100"}
\end{Highlighting}
\end{Shaded}

\subsubsection{6. Casos de uso común}\label{casos-de-uso-comuxfan}

\begin{Shaded}
\begin{Highlighting}[]
\CommentTok{\# 1. Entrada de usuario (siempre es string)}
\NormalTok{edad\_input }\OperatorTok{=} \StringTok{"25"}  \CommentTok{\# Simula input()}
\NormalTok{edad }\OperatorTok{=} \BuiltInTok{int}\NormalTok{(edad\_input)  }\CommentTok{\# Convierte a entero para operaciones}
\BuiltInTok{print}\NormalTok{(edad }\OperatorTok{+} \DecValTok{5}\NormalTok{)  }\CommentTok{\# 30}

\CommentTok{\# 2. Mostrar números en mensajes}
\NormalTok{total }\OperatorTok{=} \FloatTok{42.50}
\BuiltInTok{print}\NormalTok{(}\StringTok{"El total es: "} \OperatorTok{+} \BuiltInTok{str}\NormalTok{(total))  }\CommentTok{\# "El total es: 42.5"}

\CommentTok{\# 3. Operaciones mixtas}
\NormalTok{x }\OperatorTok{=} \StringTok{"10"}
\NormalTok{y }\OperatorTok{=} \StringTok{"5"}
\BuiltInTok{print}\NormalTok{(}\BuiltInTok{int}\NormalTok{(x) }\OperatorTok{+} \BuiltInTok{int}\NormalTok{(y))  }\CommentTok{\# 15 (convertir antes de sumar)}

\CommentTok{\# 4. Validación de datos}
\KeywordTok{def}\NormalTok{ es\_numero(texto):}
    \ControlFlowTok{try}\NormalTok{:}
        \BuiltInTok{int}\NormalTok{(texto)}
        \ControlFlowTok{return} \VariableTok{True}
    \ControlFlowTok{except} \PreprocessorTok{ValueError}\NormalTok{:}
        \ControlFlowTok{return} \VariableTok{False}
\end{Highlighting}
\end{Shaded}

\subsection{🔬 Ejemplos}\label{ejemplos}

\textbf{Mínimo ---} Conversiones básicas:

\begin{Shaded}
\begin{Highlighting}[]
\CommentTok{\# Conversiones numéricas}
\BuiltInTok{print}\NormalTok{(}\BuiltInTok{int}\NormalTok{(}\FloatTok{3.14}\NormalTok{))   }\CommentTok{\# 3}
\BuiltInTok{print}\NormalTok{(}\BuiltInTok{float}\NormalTok{(}\DecValTok{5}\NormalTok{))    }\CommentTok{\# 5.0}
\BuiltInTok{print}\NormalTok{(}\BuiltInTok{str}\NormalTok{(}\DecValTok{42}\NormalTok{))     }\CommentTok{\# "42"}
\end{Highlighting}
\end{Shaded}

\textbf{Progresivo ---} Conversiones en operaciones:

\begin{Shaded}
\begin{Highlighting}[]
\CommentTok{\# Sin casting (error)}
\CommentTok{\# print("Edad: " + 25)  \# TypeError}

\CommentTok{\# Con casting (correcto)}
\NormalTok{edad }\OperatorTok{=} \DecValTok{25}
\BuiltInTok{print}\NormalTok{(}\StringTok{"Edad: "} \OperatorTok{+} \BuiltInTok{str}\NormalTok{(edad))  }\CommentTok{\# "Edad: 25"}

\CommentTok{\# Mejor aún: f{-}strings}
\BuiltInTok{print}\NormalTok{(}\SpecialStringTok{f"Edad: }\SpecialCharTok{\{}\NormalTok{edad}\SpecialCharTok{\}}\SpecialStringTok{"}\NormalTok{)  }\CommentTok{\# "Edad: 25"}
\end{Highlighting}
\end{Shaded}

\textbf{Real ---} Sistema de registro de usuario:

\begin{Shaded}
\begin{Highlighting}[]
\CommentTok{\# Simulación de entrada (todo es texto)}
\NormalTok{nombre }\OperatorTok{=} \StringTok{"Carlos"}
\NormalTok{edad\_texto }\OperatorTok{=} \StringTok{"28"}
\NormalTok{altura\_texto }\OperatorTok{=} \StringTok{"1.75"}

\CommentTok{\# Casting para operaciones}
\NormalTok{edad }\OperatorTok{=} \BuiltInTok{int}\NormalTok{(edad\_texto)}
\NormalTok{altura }\OperatorTok{=} \BuiltInTok{float}\NormalTok{(altura\_texto)}

\CommentTok{\# Cálculos}
\NormalTok{edad\_futura }\OperatorTok{=}\NormalTok{ edad }\OperatorTok{+} \DecValTok{10}
\NormalTok{imc }\OperatorTok{=} \DecValTok{70} \OperatorTok{/}\NormalTok{ (altura }\OperatorTok{**} \DecValTok{2}\NormalTok{)  }\CommentTok{\# Peso fijo de 70kg}

\CommentTok{\# Casting para mostrar}
\BuiltInTok{print}\NormalTok{(}\StringTok{"=== Registro de Usuario ==="}\NormalTok{)}
\BuiltInTok{print}\NormalTok{(}\StringTok{"Nombre: "} \OperatorTok{+}\NormalTok{ nombre)}
\BuiltInTok{print}\NormalTok{(}\StringTok{"Edad: "} \OperatorTok{+} \BuiltInTok{str}\NormalTok{(edad) }\OperatorTok{+} \StringTok{" años"}\NormalTok{)}
\BuiltInTok{print}\NormalTok{(}\StringTok{"Altura: "} \OperatorTok{+} \BuiltInTok{str}\NormalTok{(altura) }\OperatorTok{+} \StringTok{" m"}\NormalTok{)}
\BuiltInTok{print}\NormalTok{(}\StringTok{"Edad en 10 años: "} \OperatorTok{+} \BuiltInTok{str}\NormalTok{(edad\_futura))}
\BuiltInTok{print}\NormalTok{(}\SpecialStringTok{f"IMC estimado: }\SpecialCharTok{\{}\NormalTok{imc}\SpecialCharTok{:.2f\}}\SpecialStringTok{"}\NormalTok{)  }\CommentTok{\# F{-}strings sin casting necesario}
\end{Highlighting}
\end{Shaded}

\subsection{⚠️ Errores comunes}\label{errores-comunes}

{\def\LTcaptype{none} % do not increment counter
\begin{longtable}[]{@{}
  >{\raggedright\arraybackslash}p{(\linewidth - 4\tabcolsep) * \real{0.2358}}
  >{\raggedright\arraybackslash}p{(\linewidth - 4\tabcolsep) * \real{0.3821}}
  >{\raggedright\arraybackslash}p{(\linewidth - 4\tabcolsep) * \real{0.3821}}@{}}
\toprule\noalign{}
\begin{minipage}[b]{\linewidth}\raggedright
Error
\end{minipage} & \begin{minipage}[b]{\linewidth}\raggedright
Síntoma
\end{minipage} & \begin{minipage}[b]{\linewidth}\raggedright
Solución
\end{minipage} \\
\midrule\noalign{}
\endhead
\bottomrule\noalign{}
\endlastfoot
\textbf{Texto no numérico a int} &
\texttt{ValueError:\ invalid\ literal} & Verifica que el texto sea un
número válido \\
\textbf{Texto no numérico a float} &
\texttt{ValueError:\ could\ not\ convert\ string\ to\ float} & Solo
strings numéricos (``3.14'', ``5'') \\
\textbf{Conversión de complex} &
\texttt{TypeError:\ can\textquotesingle{}t\ convert\ complex\ to\ int} &
No se puede convertir complex directamente \\
\textbf{Casting innecesario} & Código menos legible & Usa f-strings en
lugar de \texttt{str()} para mostrar \\
\textbf{Olvidar convertir input} & Suma concatenada en lugar de suma
matemática & Convierte \texttt{input()} a int/float antes de operar \\
\end{longtable}
}

\subsection{✅ Buenas prácticas}\label{buenas-pruxe1cticas}

\begin{enumerate}
\def\labelenumi{\arabic{enumi}.}
\item
  \textbf{Siempre convierte entrada de usuario}:

\begin{Shaded}
\begin{Highlighting}[]
\NormalTok{edad }\OperatorTok{=} \BuiltInTok{int}\NormalTok{(}\BuiltInTok{input}\NormalTok{(}\StringTok{"Edad: "}\NormalTok{))  }\CommentTok{\# input siempre retorna str}
\end{Highlighting}
\end{Shaded}
\item
  \textbf{Usa f-strings para mostrar datos} (evita casting manual):

\begin{Shaded}
\begin{Highlighting}[]
\CommentTok{\# ❌ Menos elegante}
\BuiltInTok{print}\NormalTok{(}\StringTok{"Edad: "} \OperatorTok{+} \BuiltInTok{str}\NormalTok{(edad) }\OperatorTok{+} \StringTok{" años"}\NormalTok{)}

\CommentTok{\# ✅ Más elegante y legible}
\BuiltInTok{print}\NormalTok{(}\SpecialStringTok{f"Edad: }\SpecialCharTok{\{}\NormalTok{edad}\SpecialCharTok{\}}\SpecialStringTok{ años"}\NormalTok{)}
\end{Highlighting}
\end{Shaded}
\item
  \textbf{Valida antes de convertir}:

\begin{Shaded}
\begin{Highlighting}[]
\NormalTok{valor }\OperatorTok{=} \BuiltInTok{input}\NormalTok{(}\StringTok{"Número: "}\NormalTok{)}
\ControlFlowTok{if}\NormalTok{ valor.isdigit():  }\CommentTok{\# Solo dígitos}
\NormalTok{    numero }\OperatorTok{=} \BuiltInTok{int}\NormalTok{(valor)}
\ControlFlowTok{else}\NormalTok{:}
    \BuiltInTok{print}\NormalTok{(}\StringTok{"No es un número válido"}\NormalTok{)}
\end{Highlighting}
\end{Shaded}
\item
  \textbf{Usa \texttt{try/except} para manejar errores de conversión}:

\begin{Shaded}
\begin{Highlighting}[]
\ControlFlowTok{try}\NormalTok{:}
\NormalTok{    edad }\OperatorTok{=} \BuiltInTok{int}\NormalTok{(}\BuiltInTok{input}\NormalTok{(}\StringTok{"Edad: "}\NormalTok{))}
\ControlFlowTok{except} \PreprocessorTok{ValueError}\NormalTok{:}
    \BuiltInTok{print}\NormalTok{(}\StringTok{"Error: debes ingresar un número"}\NormalTok{)}
\end{Highlighting}
\end{Shaded}
\item
  \textbf{Convierte lo más tarde posible}: mantén los datos en su tipo
  original hasta que necesites otro.
\end{enumerate}

\subsection{📝 Ejercicio}\label{ejercicio}

\textbf{Básico} --- Conversiones de tipos

Escribe un programa que:

\begin{enumerate}
\def\labelenumi{\arabic{enumi}.}
\tightlist
\item
  Cree las siguientes variables:

  \begin{itemize}
  \tightlist
  \item
    \texttt{numero\_texto} = ``10''
  \item
    \texttt{numero\_decimal} = 3.14
  \item
    \texttt{texto\_numerico} = ``5.5''
  \item
    \texttt{valor\_booleano} = True
  \end{itemize}
\item
  Realice las siguientes conversiones:

  \begin{itemize}
  \tightlist
  \item
    Convierta \texttt{numero\_texto} a entero
  \item
    Convierta \texttt{numero\_decimal} a entero (observa qué pasa)
  \item
    Convierta \texttt{texto\_numerico} a float
  \item
    Convierta \texttt{valor\_booleano} a entero
  \end{itemize}
\item
  Muestre los resultados y sus tipos
\end{enumerate}

\begin{center}\rule{0.5\linewidth}{0.5pt}\end{center}

\textbf{Solución (no mires hasta haberlo intentado):}

\begin{Shaded}
\begin{Highlighting}[]
\CommentTok{\# Datos originales}
\NormalTok{numero\_texto }\OperatorTok{=} \StringTok{"10"}
\NormalTok{numero\_decimal }\OperatorTok{=} \FloatTok{3.14}
\NormalTok{texto\_numerico }\OperatorTok{=} \StringTok{"5.5"}
\NormalTok{valor\_booleano }\OperatorTok{=} \VariableTok{True}

\CommentTok{\# Conversiones}
\NormalTok{entero\_desde\_texto }\OperatorTok{=} \BuiltInTok{int}\NormalTok{(numero\_texto)}
\NormalTok{entero\_desde\_decimal }\OperatorTok{=} \BuiltInTok{int}\NormalTok{(numero\_decimal)  }\CommentTok{\# Trunca a 3}
\NormalTok{float\_desde\_texto }\OperatorTok{=} \BuiltInTok{float}\NormalTok{(texto\_numerico)}
\NormalTok{entero\_desde\_booleano }\OperatorTok{=} \BuiltInTok{int}\NormalTok{(valor\_booleano)}

\CommentTok{\# Resultados}
\BuiltInTok{print}\NormalTok{(}\StringTok{"=== Conversiones ==="}\NormalTok{)}
\BuiltInTok{print}\NormalTok{(}\SpecialStringTok{f"\textquotesingle{}}\SpecialCharTok{\{}\NormalTok{numero\_texto}\SpecialCharTok{\}}\SpecialStringTok{\textquotesingle{} → }\SpecialCharTok{\{}\NormalTok{entero\_desde\_texto}\SpecialCharTok{\}}\SpecialStringTok{ (}\SpecialCharTok{\{}\BuiltInTok{type}\NormalTok{(entero\_desde\_texto)}\SpecialCharTok{\}}\SpecialStringTok{)"}\NormalTok{)}
\BuiltInTok{print}\NormalTok{(}\SpecialStringTok{f"}\SpecialCharTok{\{}\NormalTok{numero\_decimal}\SpecialCharTok{\}}\SpecialStringTok{ → }\SpecialCharTok{\{}\NormalTok{entero\_desde\_decimal}\SpecialCharTok{\}}\SpecialStringTok{ (}\SpecialCharTok{\{}\BuiltInTok{type}\NormalTok{(entero\_desde\_decimal)}\SpecialCharTok{\}}\SpecialStringTok{)"}\NormalTok{)}
\BuiltInTok{print}\NormalTok{(}\SpecialStringTok{f"\textquotesingle{}}\SpecialCharTok{\{}\NormalTok{texto\_numerico}\SpecialCharTok{\}}\SpecialStringTok{\textquotesingle{} → }\SpecialCharTok{\{}\NormalTok{float\_desde\_texto}\SpecialCharTok{\}}\SpecialStringTok{ (}\SpecialCharTok{\{}\BuiltInTok{type}\NormalTok{(float\_desde\_texto)}\SpecialCharTok{\}}\SpecialStringTok{)"}\NormalTok{)}
\BuiltInTok{print}\NormalTok{(}\SpecialStringTok{f"}\SpecialCharTok{\{}\NormalTok{valor\_booleano}\SpecialCharTok{\}}\SpecialStringTok{ → }\SpecialCharTok{\{}\NormalTok{entero\_desde\_booleano}\SpecialCharTok{\}}\SpecialStringTok{ (}\SpecialCharTok{\{}\BuiltInTok{type}\NormalTok{(entero\_desde\_booleano)}\SpecialCharTok{\}}\SpecialStringTok{)"}\NormalTok{)}

\CommentTok{\# Nota: los booleanos son subclase de int en Python}
\BuiltInTok{print}\NormalTok{(}\SpecialStringTok{f"}\CharTok{\textbackslash{}n}\SpecialStringTok{¿True es int? }\SpecialCharTok{\{}\BuiltInTok{isinstance}\NormalTok{(}\VariableTok{True}\NormalTok{, }\BuiltInTok{int}\NormalTok{)}\SpecialCharTok{\}}\SpecialStringTok{"}\NormalTok{)  }\CommentTok{\# True}
\end{Highlighting}
\end{Shaded}

\subsection{🗂️ Síntesis}\label{suxedntesis}

\begin{itemize}
\tightlist
\item
  \textbf{\texttt{int()}}: convierte a entero (trunca decimales,
  requiere números válidos).
\item
  \textbf{\texttt{float()}}: convierte a decimal (funciona con enteros y
  texto numérico).
\item
  \textbf{\texttt{str()}}: convierte a texto (funciona con casi todo).
\item
  \textbf{\texttt{bool()}}: convierte a booleano (reglas de
  truthy/falsy).
\item
  \textbf{F-strings}: alternativa moderna a \texttt{str()} para mostrar
  datos.
\item
  \textbf{Validación}: siempre verifica antes de convertir entrada de
  usuario.
\item
  \textbf{Errors comunes}: ValueError al convertir texto no numérico.
\end{itemize}

\subsection{🔗 Qué sigue}\label{quuxe9-sigue}

\begin{itemize}
\tightlist
\item
  \textbf{Relacionado:} tipos de datos, formateo de strings, validación
  de entrada
\item
  \textbf{Siguiente tema:} Strings --- aprenderás a trabajar con texto y
  todas sus operaciones y métodos.
\end{itemize}

\section{Publicaciones Similares}\label{publicaciones-similares}

Si te interesó este artículo, te recomendamos que explores otros blogs y
recursos relacionados que pueden ampliar tus conocimientos. Aquí te dejo
algunas sugerencias:

\begin{enumerate}
\def\labelenumi{\arabic{enumi}.}
\tightlist
\item
  \href{https://numerus-scriptum.netlify.app/python/2020-06-19-instalacion-de-anaconda/index.pdf}{\faIcon{file-pdf}}
  \href{https://numerus-scriptum.netlify.app/python/2020-06-19-instalacion-de-anaconda}{Instalacion
  De Anaconda}
\item
  \href{https://numerus-scriptum.netlify.app/python/2020-06-20-configurar-entorno-virtual-python-anaconda/index.pdf}{\faIcon{file-pdf}}
  \href{https://numerus-scriptum.netlify.app/python/2020-06-20-configurar-entorno-virtual-python-anaconda}{Configurar
  Entorno Virtual Python Anaconda}
\item
  \href{https://numerus-scriptum.netlify.app/python/2021-04-17-01-introducion-a-la-programacion-con-python/index.pdf}{\faIcon{file-pdf}}
  \href{https://numerus-scriptum.netlify.app/python/2021-04-17-01-introducion-a-la-programacion-con-python}{01
  Introducion A La Programacion Con Python}
\item
  \href{https://numerus-scriptum.netlify.app/python/2021-05-31-02-variables-expresiones-y-statements-con-python/index.pdf}{\faIcon{file-pdf}}
  \href{https://numerus-scriptum.netlify.app/python/2021-05-31-02-variables-expresiones-y-statements-con-python}{02
  Variables Expresiones Y Statements Con Python}
\item
  \href{https://numerus-scriptum.netlify.app/python/2021-06-07-03-objetos-de-python/index.pdf}{\faIcon{file-pdf}}
  \href{https://numerus-scriptum.netlify.app/python/2021-06-07-03-objetos-de-python}{03
  Objetos De Python}
\item
  \href{https://numerus-scriptum.netlify.app/python/2021-06-14-04-ejecucion-condicional-con-python/index.pdf}{\faIcon{file-pdf}}
  \href{https://numerus-scriptum.netlify.app/python/2021-06-14-04-ejecucion-condicional-con-python}{04
  Ejecucion Condicional Con Python}
\item
  \href{https://numerus-scriptum.netlify.app/python/2021-06-21-05-iteraciones-con-python/index.pdf}{\faIcon{file-pdf}}
  \href{https://numerus-scriptum.netlify.app/python/2021-06-21-05-iteraciones-con-python}{05
  Iteraciones Con Python}
\item
  \href{https://numerus-scriptum.netlify.app/python/2021-08-16-06-funciones-con-python/index.pdf}{\faIcon{file-pdf}}
  \href{https://numerus-scriptum.netlify.app/python/2021-08-16-06-funciones-con-python}{06
  Funciones Con Python}
\item
  \href{https://numerus-scriptum.netlify.app/python/2021-08-23-07-dataframes-con-python/index.pdf}{\faIcon{file-pdf}}
  \href{https://numerus-scriptum.netlify.app/python/2021-08-23-07-dataframes-con-python}{07
  Dataframes Con Python}
\item
  \href{https://numerus-scriptum.netlify.app/python/2021-11-29-08-prediccion-y-metrica-de-performance-con-python/index.pdf}{\faIcon{file-pdf}}
  \href{https://numerus-scriptum.netlify.app/python/2021-11-29-08-prediccion-y-metrica-de-performance-con-python}{08
  Prediccion Y Metrica De Performance Con Python}
\item
  \href{https://numerus-scriptum.netlify.app/python/2021-12-06-09-metodos-de-machine-learning-para-clasificacion-con-python/index.pdf}{\faIcon{file-pdf}}
  \href{https://numerus-scriptum.netlify.app/python/2021-12-06-09-metodos-de-machine-learning-para-clasificacion-con-python}{09
  Metodos De Machine Learning Para Clasificacion Con Python}
\item
  \href{https://numerus-scriptum.netlify.app/python/2021-12-13-10-metodos-de-machine-learning-para-regresion-con-python/index.pdf}{\faIcon{file-pdf}}
  \href{https://numerus-scriptum.netlify.app/python/2021-12-13-10-metodos-de-machine-learning-para-regresion-con-python}{10
  Metodos De Machine Learning Para Regresion Con Python}
\item
  \href{https://numerus-scriptum.netlify.app/python/2025-05-10-visualizacion-de-datos-con-python/index.pdf}{\faIcon{file-pdf}}
  \href{https://numerus-scriptum.netlify.app/python/2025-05-10-visualizacion-de-datos-con-python}{Visualizacion
  De Datos Con Python}
\end{enumerate}

Esperamos que encuentres estas publicaciones igualmente interesantes y
útiles. ¡Disfruta de la lectura!






\end{document}
